\scanpage{100}
\mspage{339}

But here the situation gets considerably more involved through the need to take account of the \deleted{\unsure{the|the}} conditions (54) (for example in the form (iv)) -- hence \inserted{that} of introducing $g_0,p$ jointly. It is then appropriate to introduce the quantities
\[
(60)\qquad \gamma=\frac{p-1}{2}\in\widehat{\mathbb Z}
\]
(N.B. since $p\in\widehat{\mathbb Z}^{\ast}$, its \unsure{class} in $\mathbb Z/2\mathbb Z$\inserted{$=\widehat{\mathbb Z}/2\widehat{\mathbb Z}$} is $1$, hence $p-1\in2\widehat{\mathbb Z}$, whence the possibility of finding $\gamma$),
\[
(61)\qquad h_0=g_0^{-1}\sigma_\infty(g_0)
\]

\margin{$h_0$, relative to the automorphism $\sigma_\infty$ of order $2$, \unsure{plays} a role symmetric \unsure{to} $\xi$ relative to the automorphism $\rho$ of order $3$ -- this does not show up \unsure{either} in the notation.}

\deleted{(we change our notation, previously we had taken $h_0=\sigma_\infty(g_0)^{-1}g_0$, so we take the inverse \unsure{with respect|p.v.} to the previous notation).} We check \unsure{indeed|thus} that $g$ is expressed by means of $g_0$ and of $\gamma$ (hence of $p$) by means of the formula (\uline{equivalent} to (54) (iv) or (v))
\[
(62)\quad
\boxed{g=\sigma_\infty(g_0)\ell_1^\gamma\rho(g_0)^{-1}}
\left(=(g_0\ell_0^{-\gamma})\left(\underbrace{\ell_0^\gamma h_0}_{\xi}\right)\rho(g_0\ell_0^{-\gamma})^{-1}\right)
\]

\margin{\uline{N.B.} This expression for $g$ is not \unsure{changed} if $g_0\longmapsto g_0\ell_0^\alpha$.}

Thus the $u$ that interest us are described by the pairs $(g_0,\gamma)$ $(g_0\in\widehat{\Pi}_{0,3}/L_0,\ \gamma\in\widehat{\mathbb Z})$ \unsure{satisfying} the conditions
\[
(63)\qquad 2\gamma+1\in\widehat{\mathbb Z}^{\ast}\qquad (\mathrm{NB}\quad 2\gamma+1=p=\chi(u))
\]
and condition (51), where $g$ is expressed as a function of $g_0$ and of $\gamma$ by (62). This \deleted{formula}\inserted{relation} (51) is in turn equivalent to the relation
\[
(64)\qquad\boxed{\xi\rho(\xi)\rho^2(\xi)=1}
\]
where
\[
(65)\quad\xi=\ell_0^\gamma h_0
\qquad\unsure{\deleted{\text{with}}|\text{with}}\ (66)\quad
\begin{aligned}
h_0&\overset{\mathrm{dfn}}{=}
g_0^{-1}\sigma_\infty(g_0)\\
&{}=[g_0^{-1},\sigma_\infty]
\end{aligned}
\]

\margin{\uline{N.B.} We \unsure{call} $h_0$ here what we previously denoted by \unsure{$h_0^{-1}$}.}

\scanpage{101}

Thus, the $u$ are \unsure{parametrized} by the pairs
$(g_0\in\widehat{\Pi}_{0,3}/L_0,\ \gamma\in\widehat{\mathbb Z})$, satisfying conditions
(63) and (64), which (by means of (65)) is written in explicit form
\[
\tag{\unsure{65bis}}
\boxed{
\bigl(\ell_0^\gamma\underbrace{g_0^{-1}\sigma_\infty(g_0)}_{h_0}\bigr)
\rho\bigl(\ell_0^\gamma\underbrace{g_0^{-1}\sigma_\infty(g_0)}_{h_0}\bigr)
\rho^2\bigl(\ell_0^\gamma\underbrace{g_0^{-1}\sigma_\infty(g_0)}_{h_0}\bigr)=1
}
\]
\margin{it would be better to regard this equation as an ``equation in $h_0$'' \unsure{playing} a role symmetric to ``the equation in $\xi$'' (68)}
(``equation in $g_0$'' --- as opposed to the ``equation in $\xi$'' below).
(Moreover, we \unsure{must} add the preliminary condition, for the record...)

We can also consider that \deleted{$g$ is} $u$ (via $g$) is defined by
$(\xi\in\widehat{\Pi}_{0,3},\ \gamma\in\widehat{\mathbb Z})$, which indeed allow us to recover
\[
h_0=\ell_0^{-\gamma}\xi
\]
by (65), then $g_0$ by (66)
\[
g_0^{-1}\sigma_\infty(g_0)=h_0
\]
which indeed determines $g_0$ uniquely, thanks to the fact (which we admit provisionally)
that $\widehat{\Pi}_{0,3}^{\sigma_\infty}=\{1\}$.
The possibility of finding $g_0$ satisfying this condition is equivalent (we shall admit this here too) to the condition
\[
\tag{67}
\sigma_\infty(h_0)h_0=1
\]
which, in terms of $\xi,\gamma$, is also written in the \unsure{form}

\scanpage{102}
\mspage{340}

\[
\tag{68}
\boxed{\sigma_\infty(\xi)=\ell_1^\gamma\xi^{-1}\ell_0^\gamma}.
\]

Equations (64), (68) are called the ``equations in $\xi$'' --- as opposed to the ``equation in $g_0$'' (66), which at first sight has the advantage of being a single one, and the disadvantage of being more \unsure{complicated}.
\margin{rather ``the equation in $h_0$''}
In summary, the $u$ that interest us
$(u\in\mathrm{Aut}_{\mathrm{lac}}(\widehat{\mathfrak S}_{0,3}^{+},\sigma_\infty)')$
are \unsure{given|described} by the pairs
\deleted{$\unsure{(\xi,\gamma)}$}\inserted{$(\xi\in\widehat{\Pi}_{0,3}^{+},\ \gamma\in\widehat{\mathbb Z})$}, satisfying the conditions
\[
\begin{aligned}
\text{a)}\quad&2\gamma+1\in\widehat{\mathbb Z}^{*} &&(63)\\
\text{b)}\quad&\xi\rho(\xi)\rho^2(\xi)=1 &&(64)\\
\text{c)}\quad&\sigma_\infty(\xi)=\ell_1^\gamma\xi^{-1}\ell_0^\gamma &&(68)
\end{aligned}
\]
d) Preliminary condition for the record (53)
(which apparently are not \unsure{expressed} directly in a convenient way in terms of $\xi,\gamma$).

For two pairs $(\xi,\unsure{\gamma})$, $(\xi',\unsure{\gamma}')$ to define the same $u$, it is necessary and sufficient that we have
\[
\tag{69}
\unsure{\gamma}=\unsure{\gamma}',\qquad
\xi'=\ell_0^{-\alpha}\xi \ell_1^\alpha\qquad(\alpha\in\widehat{\mathbb Z}).
\]

Condition (64) can also be \unsure{replaced} by a parametric representation
\[
\tag{70}
\left\{
\begin{aligned}
\xi&=\lambda\rho(\lambda)^{-1}\qquad\text{or}\\
\xi&=\lambda \ell_0^{-1}\rho(\lambda)^{-1}
\end{aligned}
\right.
\]
with $\lambda\in\widehat{\Pi}_{0,3}$ --- in which case equation (68)

\scanpage{103}

can be interpreted as an ``equation in $\lambda$'', which according to the two cases takes the form
\[
\tag{71}
\left\{
\begin{aligned}
\sigma_\infty(\lambda\rho(\lambda)^{-1})
 &=\ell_1^\gamma(\rho(\lambda)\lambda^{-1})\ell_0^\gamma\\
&\text{or}\\
\sigma_\infty(\lambda\unsure{\ell_0^{-1}}\rho(\lambda)^{-1})
 &=\ell_1^\gamma(\rho(\lambda)\ell_0\lambda^{-1})\ell_0^\gamma,
\end{aligned}
\right.
\]

Here for $(\lambda,\gamma)$ and $(\lambda',\gamma')$ to define the same $u$, it is necessary and sufficient that we have \deleted{\illegible{1}}
\[
\tag{72}
\lambda'=\ell_0^{-\alpha}\lambda,\qquad\gamma=\gamma'.
\]

\uline{Remarks} 1) Up to the ``harmless'' factor $\ell_0^\gamma$, \unsure{which|that} distinguishes $h_0$ and $\xi$ (65), we can say that the \unsure{search} for the
$u\in\mathrm{Aut}_{\mathrm{lac}}(\widehat{\mathfrak S}_{0,3}^{+},\sigma_\infty)$
is equivalent to that for the \deleted{\unsure{equations}} \uline{simultaneous} solutions of the two cocycle equations
\[
\tag{73}
\left\{
\begin{aligned}
\xi\rho(\xi)\rho^2(\xi)&=1\\
h_0\sigma_\infty(h_0)&=1
\end{aligned}
\right.
\]
relative \deleted{respectively} to the groups $\mathbb Z/3\mathbb Z$, operating on $\widehat{\Pi}_{0,3}$ via $\rho$, and $\mathbb Z/2\mathbb Z$, operating via $\sigma_\infty$ --- whose solutions in parametric forms are respectively
\[
\tag{74}
\left\{
\begin{aligned}
\xi&=\lambda\rho(\lambda)^{-1}\quad\text{or}\quad
 \xi=\lambda \ell_0^{-1}\rho(\lambda)^{-1}\\
&\quad\text{(according as }\gamma\equiv0\ (3)\text{ or }\gamma\equiv\unsure{1|-1}\ (3)\\
&\qquad\text{i.e. }p\equiv1\ (3)\text{ or }p\equiv-1\ (3)\text{)}\\
h_0&\unsure{=|\equiv}\unsure{g_0^{-1}\sigma_\infty(g_0)}
\end{aligned}
\right.
\]
\margin{N.B. expression (62) for $g$ brings out equally well \unsure{this double} \unsure{intervention} of $\sigma_\infty$ and of $\rho$ in the determination of the $u$!!!}
\margin{(N.B. we could put \unsure{everywhere} $g_0$ for $g_0^{-1}$...)}

\scanpage{104}
\mspage{341}

3) The relation between the $\lambda_0$ in the parametric representation of $g$ (58), and the $\lambda$ in that of $\xi$ (74 a) is implicit in the last expression (62) for $g$, namely
\[
\tag{75}
\lambda_0=g_0\ell_0^{-\gamma}\lambda.
\]

4) Let us note here also, that the quantities
\[
\begin{array}{cccccccc}
p,&\gamma,&g,&g_0,&h_0,&\xi,&\lambda_0,&\lambda\\
\rotatebox{-90}{$\in$}&\rotatebox{-90}{$\in$}&\rotatebox{-90}{$\in$}&\rotatebox{-90}{$\in$}&\rotatebox{-90}{$\in$}&\rotatebox{-90}{$\in$}&\rotatebox{-90}{$\in$}&\rotatebox{-90}{$\in$}\\
\widehat{\mathbb Z}^{*}&\widehat{\mathbb Z}&\widehat{\Pi}_{0,3}&\widehat{\Pi}_{0,3}/L_0&&&\widehat{\Pi}_{0,3}&
\end{array}
\]
\inserted{resp. $\lambda'_0,\lambda'$}
(with \inserted{for $u$ constant} the indeterminacies specific to $g_0,h_0,\xi,\lambda$) can be considered \inserted{for $u$ variable,} as functions on the group
$\mathrm{Autext}_{\mathrm{lac}}(\widehat{\mathfrak S}_{0,3}^{+})'\simeq\mathrm{Autext}_{\mathrm{lac}}(\widehat{\Pi}_{0,3})'$
(subgroup of the outer automorphisms \unsure{of the group} with loops $\widehat{\Pi}_{0,3}$ that commute with $\mathfrak S_3$), with values in the sets
$\widehat{\mathbb Z}^{*},\widehat{\mathbb Z},\widehat{\Pi}_{0,3},\widehat{\Pi}_{0,3}/L_0,\ ,\ ,\widehat{\Pi}_{0,3},\ $.
In particular, we thus find 8 functions \inserted{\unsure{continuous}} on $\Gamma_{\mathbb Q}$, with values in the appropriate spaces --- starting with the cyclotomic character.
\margin{We shall see (\illegible{2} \unsure{less}!) that we can \uline{normalize} $g_0$, and thereby also $h_0,\xi,\lambda,\lambda_0$, so as to remove all ambiguity in \unsure{its|these} \unsure{definition|definitions}...}

\scanpage{105}

Let us note in this connection the equivalent conditions
\[
\tag{75bis}
\left\{
\begin{aligned}
\text{(i)}\quad &p=1\\
\text{(ii)}\quad &\gamma=0\\
\text{(iii)}\quad &g=\sigma_\infty(g_0)\rho(g_0)^{-1}\\
\text{(iv)}\quad &h_0=\xi\\
\text{(v)}\quad &\sigma_\infty(\xi)=\xi^{-1}
\qquad\text{(``equation in $\xi$'' simplified!)}\\
\text{(vi)}\quad &h_0\rho(h_0)\rho^2(h_0)=1
\end{aligned}
\right.
\]
all meaning that $\chi(u)=1$. Thus the search for the
$u\in\mathrm{Aut}_{\mathrm{lac}}(\widehat{\mathfrak S}_{0,3}^{+},\sigma_\infty)'$
\uline{with multiplier 1}
($u\in\mathrm{Aut}_{\mathrm{lac}}(\widehat{\mathfrak S}_{0,3}^{+},\sigma_\infty)^{\prime0}$)
is thus indeed equivalent to that for the simultaneous solutions of the equations
\inserted{of cocycles} (73), or of the parametric equations (``coboundary'') (74).

2) Returning to the general case (arbitrary multiplier), the only case that I know how to write down explicitly
\inserted{(except for $u=1$)} \unsure{is} the one that corresponds to
$\tau_\infty\in\unsure{\mathfrak S_{0,3}^{*}}$,
which corresponds to the characteristics
\[
\tag{75ter}
\left\{
\begin{aligned}
&u=\tau_\infty,\quad\chi(u)=p=-1,\quad\gamma=-1\\
&g=\ell_1^{-1}\\
&g_0=1,\quad h_0=1,\quad\xi=\ell_0^{-1}\\
&\lambda'_0=1,\quad\lambda'=\unsure{\ell_0^\gamma g_0^{-1}\lambda_0}=\ell_0^{-1}\\
&\inserted{[\,\lambda_0=\tau_\infty,\quad\lambda=\ell_0^{-1}\tau_\infty\,]}
\end{aligned}
\right.
\]
\margin{$\leftarrow$ cf. below}

It would certainly be no luxury \unsure{to manage} to write down a bunch of other $u$, to see in particular whether the \unsure{crowds|crowd} of necessary conditions that we are led to derive for the elements coming from $\unsure{\Gamma_{\mathbb Q}}$, are not automatically satisfied!